\documentclass[11pt]{article}

% ===================== PACKAGES =====================
\usepackage[a4paper,margin=1in]{geometry}
\usepackage{amsmath,amssymb}
\usepackage{enumitem}
\usepackage{fancyhdr}
\usepackage{xcolor} 
\usepackage{algorithm}
\usepackage{algpseudocode}

% ===================== HEADER & FOOTER =====================
\pagestyle{fancy}
\fancyhf{}
\lhead{CSE 400: Fundamentals of Probability in Computing}
\rhead{Tutorial-2 Scribe}
\cfoot{\thepage}

% ===================== TITLE =====================
\title{
    \normalsize School of Engineering and Applied Science (SEAS), Ahmedabad University \\
    \vspace{0.2cm}
    \textbf{CSE 400: Fundamentals of Probability in Computing}\\
    \Large Tutorial-2 Lecture Scribe
}
\author{} 
\date{}

\begin{document}
\maketitle

\vspace{-1cm}
\begin{center}
    \begin{tabular}{ll}
        \textbf{Date:} & {February 25, 2026 \hspace{2.2in}} \\ [1.5ex]
        
    \end{tabular}
\end{center}

\hrule
\vspace{0.5cm}

% ===================== CONTENT =====================

\section*{Question 1}
\textbf{Problem Statement:} A point is chosen at random on a line segment of length $L$. Interpret this statement and find the probability that the ratio of the shorter to the longer segment is less than $\frac{1}{4}$.

\subsection*{Step 1: Interpretation}
Choosing a point at random on a line segment of length $L$ means:
\[
X \sim \text{Uniform}(0, L)
\]
where $X$ represents the distance of the chosen point from one endpoint.
The segment is divided into two parts: Length $X$ and Length $L - X$.

The shorter segment is:
\[
\min(X, L - X)
\]
The longer segment is:
\[
\max(X, L - X)
\]
We need:
\[
\frac{\min(X, L - X)}{\max(X, L - X)} < \frac{1}{4}
\]

\subsection*{Step 2: Case Analysis}
\textbf{Case 1: $X \le L - X$} \\
This occurs when:
\[
X \le \frac{L}{2}
\]
Then:
\[
\text{Shorter} = X, \quad \text{Longer} = L - X
\]
Condition:
\[
\frac{X}{L - X} < \frac{1}{4}
\]
Multiply both sides by $(L - X > 0)$:
\[
X < \frac{L - X}{4}
\]
Multiply both sides by 4:
\begin{align*}
4X &< L - X \\
5X &< L \\
X &< \frac{L}{5}
\end{align*}

\textbf{Case 2: $X > \frac{L}{2}$} \\
Then:
\[
\text{Shorter} = L - X, \quad \text{Longer} = X
\]
Condition:
\[
\frac{L - X}{X} < \frac{1}{4}
\]
Multiply both sides by $(X > 0)$:
\[
L - X < \frac{X}{4}
\]
Multiply both sides by 4:
\begin{align*}
4L - 4X &< X \\
4L &< 5X \\
X &> \frac{4L}{5}
\end{align*}

\subsection*{Step 3: Probability Computation}
Favorable region:
\[
X < \frac{L}{5} \quad \text{or} \quad X > \frac{4L}{5}
\]
Total length of favorable region:
\[
\frac{L}{5} + \frac{L}{5} = \frac{2L}{5}
\]
Since total sample space length is $L$,
\[
P = \frac{2L/5}{L} = \frac{2}{5}
\]

\section*{Question 2}
(White region: $\mathcal{N}(4,4)$, Black region: $\mathcal{N}(6,9)$, Black fraction = $\alpha$) \\
We are given reading $x = 5$.



\subsection*{Step 1: Densities at $x=5$}
White density:
\begin{align*}
f_W(5) &= \frac{1}{\sqrt{2\pi(4)}} \exp\left(-\frac{(5-4)^2}{2(4)}\right) \\
&= \frac{1}{2\sqrt{2\pi}} \exp\left(-\frac{1}{8}\right)
\end{align*}

Black density:
\begin{align*}
f_B(5) &= \frac{1}{\sqrt{2\pi(9)}} \exp\left(-\frac{(5-6)^2}{2(9)}\right) \\
&= \frac{1}{3\sqrt{2\pi}} \exp\left(-\frac{1}{18}\right)
\end{align*}

\subsection*{Step 2: Equal Error Condition}
Probability of misclassifying white:
\[
(1-\alpha) f_W(5)
\]
Probability of misclassifying black:
\[
\alpha f_B(5)
\]
Set equal:
\[
(1-\alpha) f_W(5) = \alpha f_B(5)
\]
Substitute expressions:
\[
(1-\alpha)\frac{1}{2} e^{-1/8} = \alpha \frac{1}{3} e^{-1/18}
\]
Solve for $\alpha$:
\[
\alpha = \frac{\frac{1}{2}e^{-1/8}}{\frac{1}{2}e^{-1/8} + \frac{1}{3}e^{-1/18}}
\]
This is the required value of $\alpha$.

\section*{Question 3}
\subsection*{(a) $X \sim \text{Uniform}(0,A)$}
We minimize:
\[
E|X-a| = \int_0^A |x-a| \frac{1}{A} \, dx
\]
Split at $a$:
\[
= \frac{1}{A} \left[ \int_0^a (a-x) \, dx + \int_a^A (x-a) \, dx \right]
\]
Compute integrals:
\begin{align*}
\int_0^a (a-x) \, dx &= \left( ax - \frac{x^2}{2} \right)\Bigg|_0^a = a^2 - \frac{a^2}{2} = \frac{a^2}{2} \\
\int_a^A (x-a) \, dx &= \left( \frac{x^2}{2} - ax \right)\Bigg|_a^A = \left(\frac{A^2}{2} - aA\right) - \left(\frac{a^2}{2}-a^2\right) = \frac{A^2}{2} - aA + \frac{a^2}{2}
\end{align*}
Sum:
\[
E|X-a| = \frac{1}{A} \left( a^2 - aA + \frac{A^2}{2} \right)
\]
Differentiate with respect to $a$:
\[
\frac{d}{da} = \frac{1}{A}(2a - A)
\]
Set to zero:
\begin{align*}
2a - A &= 0 \\
a &= \frac{A}{2}
\end{align*}

\subsection*{(b) $X \sim \text{Exponential}(\lambda)$}
\[
E|X-a| = \int_0^\infty |x-a| \lambda e^{-\lambda x} \, dx
\]
Split at $a$. Differentiate w.r.t. $a$ and set equal to zero.
This yields the median condition:
\[
P(X \le a) = \frac{1}{2}
\]
For exponential:
\begin{align*}
1 - e^{-\lambda a} &= \frac{1}{2} \\
e^{-\lambda a} &= \frac{1}{2} \\
a &= \frac{\ln 2}{\lambda}
\end{align*}

\section*{Question 4}
\textbf{Mixture of exponentials.} \\
We compute:
\[
P(X > t+s \mid X > t) = \frac{P(X>t+s)}{P(X>t)}
\]
Compute numerator:
\[
P(X>t+s) = p_1 e^{-\lambda_1(t+s)} + p_2 e^{-\lambda_2(t+s)}
\]
Denominator:
\[
P(X>t) = p_1 e^{-\lambda_1 t} + p_2 e^{-\lambda_2 t}
\]
Thus:
\[
\frac{p_1 e^{-\lambda_1(t+s)} + p_2 e^{-\lambda_2(t+s)}}{p_1 e^{-\lambda_1 t} + p_2 e^{-\lambda_2 t}}
\]

\section*{Question 5}
Given $X \sim \text{Poisson}(\lambda_1)$, $Y \sim \text{Poisson}(\lambda_2)$. \\
We want:
\[
P(X=k \mid X+Y=n)
\]
Using definition:
\[
= \frac{P(X=k,Y=n-k)}{P(X+Y=n)}
\]
Numerator:
\[
\frac{e^{-\lambda_1}\lambda_1^k}{k!} \cdot \frac{e^{-\lambda_2}\lambda_2^{n-k}}{(n-k)!}
\]
Denominator:
\[
\frac{e^{-(\lambda_1+\lambda_2)}(\lambda_1+\lambda_2)^n}{n!}
\]
Simplify:
\[
\binom{n}{k} \left(\frac{\lambda_1}{\lambda_1+\lambda_2}\right)^k \left(\frac{\lambda_2}{\lambda_1+\lambda_2}\right)^{n-k}
\]

\section*{Question 6}
\textbf{Transformation} $Y=g(X)$, $X \sim \text{Uniform}[-2,2]$. \\
CDF:
\[
F_Y(y) = P(g(X)\le y)
\]
Find preimages and compute probabilities using:
\[
P(a<X<b) = \frac{b-a}{4}
\]
Differentiate:
\[
f_Y(y) = \frac{d}{dy}F_Y(y)
\]
\textit{(Graph construction follows from computed piecewise expressions.)}

\section*{Question 7}
\textbf{Binomial model:} \\
Number of up moves $S \sim \text{Bin}(1000,0.52)$. \\
Need:
\[
u^k d^{1000-k} \ge 1.3
\]
Take log:
\[
k\ln u + (1000-k)\ln d \ge \ln 1.3
\]
Solve for $k$. \\
Approximate using normal:
\begin{align*}
\mu &= 1000(0.52) = 520 \\
\sigma^2 &= 1000(0.52)(0.48)
\end{align*}
Standardize and compute probability.

\section*{Question 8}
Given $T \sim \text{Exponential}(\lambda=1/2)$.

\subsection*{(a)}
\[
P(T>2) = e^{-(1/2)(2)} = e^{-1}
\]

\subsection*{(b)}
\[
P(T>10 \mid T>9) = \frac{e^{-5}}{e^{-4.5}} = e^{-0.5}
\]

\section*{Question 9}
\[
Z = \sqrt{X^2+Y^2}
\]



Use 2D CDF method:
\[
F_Z(z) = \iint_{x^2+y^2 \le z^2} f_X(x)f_Y(y) \, dx \, dy
\]
Differentiate:
\[
f_Z(z) = \int_0^{2\pi} f_X(z\cos\theta) f_Y(z\sin\theta) z \, d\theta
\]

\section*{Question 10}
Let $M_n = \max(X_1,\dots,X_n)$.
\[
P(M_n \le x) = F(x)^n
\]
Joint probability:
\[
P(M_n \le x, M_{n+1} \le y)
\]
For $x \le y$:
\[
F(x)^n F(y)
\]
For $x > y$:
\[
F(y)^{n+1}
\]

\bigskip
\hrule
\vspace{0.2cm}
\begin{center}
    \small \textit{End of Complete Tutorial-2 Scribe}
\end{center}

\end{document}